\chapter{多项式}

本章讨论域 $\F$ 上的一元多项式环 $\F[x]$。
在线性代数中，多项式并不只是代数练习对象；真正关键的是，
若 $T\in\End(V)$，则每个多项式
\[
p(x)=a_0+a_1x+\cdots+a_mx^m
\]
都可以作用在算子上，定义为
\[
p(T)=a_0\Id+a_1T+\cdots+a_mT^m.
\]
后续关于特征值、上三角化与谱定理的论证，
都要依赖本章建立的多项式工具。

\section{多项式环 $\F[x]$}

\begin{definition}
域 $\F$ 上的\emph{多项式}是形如
\[
a_0+a_1x+\cdots+a_nx^n
\]
的表达式，其中 $n\in\N$，$a_0,\dots,a_n\in\F$。
所有这样的多项式构成集合 $\F[x]$。
\end{definition}

\begin{definition}
若 $p(x)=a_0+a_1x+\cdots+a_nx^n$ 且 $a_n\neq 0$，
则称 $n$ 为 $p$ 的\emph{次数}，记作 $\deg p=n$。
零多项式的次数不定义。
\end{definition}

多项式相等是指对应幂次的系数逐项相等。
加法与乘法按通常方式定义。
$\F[x]$ 在这些运算下是交换环，而且没有零因子。

\begin{proposition}
设 $p,q\in\F[x]$ 且都非零。
则
\[
\deg(pq)=\deg p+\deg q.
\]
此外，若 $\deg p\neq \deg q$，则
\[
\deg(p+q)=\max\{\deg p,\deg q\}.
\]
\end{proposition}

\begin{proof}
设 $\deg p=m$，$\deg q=n$，其首项系数分别为 $a_m,b_n$。
则 $pq$ 的 $x^{m+n}$ 项系数为 $a_mb_n\neq 0$，
故 $\deg(pq)=m+n$。

若 $m\neq n$，不妨设 $m>n$。
则 $p+q$ 的 $x^m$ 项系数仍为 $p$ 的首项系数，非零，
故 $\deg(p+q)=m$。
\end{proof}

\begin{remark}
当 $\deg p=\deg q$ 时，$p+q$ 的最高次项可能抵消，
因此 $\deg(p+q)$ 可能小于 $\deg p$。
这正是多项式代数中最常见的细节之一。
\end{remark}

\begin{example}
在 $\R[x]$ 中，
\[
(x^2+1)+( -x^2+3x)=3x+1,
\]
所以两个二次多项式之和可以降为一次多项式。
\end{example}

\begin{example}
在 $\F$ 上，若 $p(x)=a_0+a_1x+\cdots+a_nx^n$，则对任意 $\lambda\in\F$，
定义
\[
p(\lambda)=a_0+a_1\lambda+\cdots+a_n\lambda^n.
\]
这称为将多项式在 $\lambda$ 处取值。
\end{example}

\section{整除与因式分解}

\begin{definition}
设 $p,q\in\F[x]$。
若存在 $s\in\F[x]$ 使得 $q=ps$，
则称 $p$ \emph{整除} $q$，记作 $p\mid q$。
\end{definition}

\begin{theorem}[带余除法]
设 $p,q\in\F[x]$ 且 $q\neq 0$。
则存在唯一的多项式 $s,r\in\F[x]$，使得
\[
p=qs+r,
\]
并且满足 $r=0$ 或 $\deg r<\deg q$。
\end{theorem}

\begin{proof}
若 $p=0$，取 $s=r=0$ 即可。
现设 $p\neq 0$。
若 $\deg p<\deg q$，取 $s=0,r=p$。

若 $\deg p\ge \deg q$，设 $\deg p=m$，$\deg q=n$，
首项系数分别为 $a,b$。
令
\[
t(x)=\frac{a}{b}x^{m-n}.
\]
则 $p-tq$ 的最高次项被消去，所以 $\deg(p-tq)<m$。
对次数作归纳，存在 $u,r$ 使
\[
p-tq=qu+r,
\]
其中 $r=0$ 或 $\deg r<\deg q$。
于是
\[
p=q(t+u)+r.
\]
存在性得证。

若还有 $p=qs_1+r_1=qs_2+r_2$，且 $\deg r_1,\deg r_2<\deg q$，
则
\[
q(s_1-s_2)=r_2-r_1.
\]
若 $s_1\neq s_2$，左边次数至少为 $\deg q$，右边次数却小于 $\deg q$，矛盾。
故 $s_1=s_2$，进而 $r_1=r_2$。
\end{proof}

\begin{definition}
设 $p_1,\dots,p_k\in\F[x]$ 不全为零。
若多项式 $d$ 同时整除每个 $p_j$，并且任一公共因子都整除 $d$，
则称 $d$ 是它们的\emph{最大公因子}。
通常约定取首项系数为 $1$ 的那个，称为首一最大公因子。
\end{definition}

\begin{proposition}
任意两个不全为零的多项式 $p,q\in\F[x]$ 都有唯一的首一最大公因子。
\end{proposition}

\begin{proof}
利用带余除法作欧几里得算法：
\[
p=q s_1+r_1,\quad q=r_1s_2+r_2,\quad \dots
\]
每一步余式次数严格下降，因此过程必在有限步后结束，
最后一个非零余式记为 $d$。
像整数情形一样，由回代可知 $d$ 是 $p,q$ 的线性组合，
因此任何公共因子都整除 $d$；另一方面 $d$ 整除前一项，也整除再前一项，
最终整除 $p,q$。
故 $d$ 为最大公因子。
乘以非零常数可将它改为首一形式，且此形式唯一。
\end{proof}

\begin{proposition}
若 $\gcd(p,q)=1$ 且 $p\mid qr$，则 $p\mid r$。
\end{proposition}

\begin{proof}
因为 $\gcd(p,q)=1$，由欧几里得算法的回代可得存在 $a,b\in\F[x]$ 使
\[
ap+bq=1.
\]
两边乘以 $r$ 得
\[
ar p+bqr=r.
\]
右边两项都被 $p$ 整除，因此 $p\mid r$。
\end{proof}

\section{零点}

\begin{definition}
设 $p\in\F[x]$，$\lambda\in\F$。
若 $p(\lambda)=0$，则称 $\lambda$ 是 $p$ 的\emph{零点}或\emph{根}。
\end{definition}

\begin{theorem}[因子定理]
设 $p\in\F[x]$，$\lambda\in\F$。
则
\[
p(\lambda)=0 \iff (x-\lambda)\mid p(x).
\]
\end{theorem}

\begin{proof}
将 $p$ 除以 $x-\lambda$，可得
\[
p(x)=(x-\lambda)q(x)+r
\]
其中 $r\in\F$ 为常数。
代入 $x=\lambda$ 得
\[
p(\lambda)=r.
\]
因此 $p(\lambda)=0$ 当且仅当 $r=0$，
也即当且仅当 $(x-\lambda)\mid p(x)$。
\end{proof}

\begin{corollary}
非零多项式的零点个数不超过其次数。
\end{corollary}

\begin{proof}
设 $\deg p=n$。
若 $\lambda_1,\dots,\lambda_k$ 是互异零点，
则由因子定理，
\[
(x-\lambda_1)\cdots(x-\lambda_k)\mid p(x).
\]
左边次数为 $k$，故 $k\le n$。
\end{proof}

\begin{definition}
若 $(x-\lambda)^m\mid p$ 且 $(x-\lambda)^{m+1}\nmid p$，
则称 $\lambda$ 是 $p$ 的\emph{重数为 $m$ 的零点}。
\end{definition}

\begin{remark}
本书不依赖导数来讨论重根。
对我们而言，重数的本质是线性因子在分解中重复出现的次数。
\end{remark}

\begin{example}
在 $\R[x]$ 中，
\[
(x-2)^3(x+1)
\]
有零点 $2$，其重数为 $3$；零点 $-1$，其重数为 $1$。
\end{example}

\section{$\C$ 上的分解}

\begin{theorem}[代数基本定理，叙述]
每个非常数复系数多项式在 $\C$ 中至少有一个零点。
\end{theorem}

我们不证明该定理，只使用其结论。
它与线性代数结合后，会导致复向量空间上的算子一定存在特征值。

\begin{corollary}
每个非常数 $p\in\C[x]$ 都可分解为一次因子的乘积：
\[
p(x)=c(x-\lambda_1)\cdots(x-\lambda_n),
\]
其中 $c\in\C\setminus\{0\}$。
\end{corollary}

\begin{proof}
对 $\deg p$ 归纳。
次数为 $1$ 时显然成立。
若 $\deg p=n\ge 2$，由代数基本定理存在 $\lambda\in\C$ 使 $p(\lambda)=0$。
由因子定理，$p=(x-\lambda)q$，其中 $\deg q=n-1$。
对 $q$ 用归纳假设即可。
\end{proof}

\begin{corollary}
$\C[x]$ 中不可约多项式恰为一次多项式。
\end{corollary}

\begin{proof}
一次多项式显然不可约。
若 $p\in\C[x]$ 的次数至少为 $2$，
则上面的分解表明它可写成两个次数正的多项式之积，故可约。
\end{proof}

\section{$\R$ 上的分解}

\begin{proposition}
若 $p\in\R[x]$ 且 $z\in\C$ 是其非实零点，
则共轭数 $\overline z$ 也是零点。
\end{proposition}

\begin{proof}
设
\[
p(x)=a_0+a_1x+\cdots+a_nx^n
\quad (a_j\in\R).
\]
若 $p(z)=0$，则取共轭得
\[
0=\overline{p(z)}=a_0+a_1\overline z+\cdots+a_n\overline z^n=p(\overline z).
\]
\end{proof}

\begin{theorem}
每个非常数 $p\in\R[x]$ 都可分解为一次或二次实系数不可约因子的乘积。
\end{theorem}

\begin{proof}
把 $p$ 看作复系数多项式，在 $\C[x]$ 中完全分解为一次因子。
其中每个实零点对应一次实因子 $(x-\lambda)$。
每对非实共轭零点 $a\pm bi$ 对应乘积
\[
(x-(a+bi))(x-(a-bi))=(x-a)^2+b^2,
\]
这是实系数二次多项式。
把所有因子配对，即得所求分解。
\end{proof}

\begin{corollary}
$\R[x]$ 中不可约多项式恰为一次多项式和没有实零点的二次多项式。
\end{corollary}

\begin{proof}
一次多项式不可约。
若二次多项式有实零点，则由因子定理可分解为两个一次因子；
故没有实零点的二次多项式不可约。
若次数至少为 $3$，由上一定理可分解为低次因子，因此可约。
\end{proof}

\begin{remark}
这里的“没有实零点的二次多项式”可视作实数域上不可再分解的基本块。
在后面研究实向量空间上的算子时，这些二次因子会对应二维不变子空间。
\end{remark}

\section{多项式作用于算子}

本章真正服务于线性代数的核心思想是：
多项式可以代入算子。

\begin{definition}
设 $V$ 是向量空间，$T\in\End(V)$，
若
\[
p(x)=a_0+a_1x+\cdots+a_nx^n\in\F[x],
\]
则定义
\[
p(T)=a_0\Id+a_1T+\cdots+a_nT^n\in\End(V).
\]
\end{definition}

\begin{proposition}
对任意 $p,q\in\F[x]$ 与 $T\in\End(V)$，有
\[
(p+q)(T)=p(T)+q(T),\qquad (pq)(T)=p(T)q(T).
\]
\end{proposition}

\begin{proof}
第一式由线性直接得到。
第二式来自算子乘法与多项式乘法的定义：
把 $p(T)$ 与 $q(T)$ 展开后，按同次幂合并即可。
其中关键是 $T^jT^k=T^{j+k}$，并且各个 $T^m$ 两两可交换，
因为它们都由同一个算子 $T$ 的幂构成。
\end{proof}

\begin{remark}
后面若出现 $p(T)=0$，我们就可以把 $p$ 在 $\C[x]$ 或 $\R[x]$ 中分解，
从而把关于算子的一个高次关系拆成若干低次关系。
这正是 Axler 风格论证的基本技术。
\end{remark}

\section*{习题}

\subsection*{基础题}
\begin{enumerate}[label=\textbf{\arabic*.}, leftmargin=*, itemsep=0.5em]
\item \basic 设 $p(x)=3-2x+x^3$，写出它的次数与各次幂的系数。
\item \basic 计算 $(2+x)+(1-3x+4x^2)$。
\item \basic 计算 $(1+x)(1-x+x^2)$。
\item \basic 求 $\deg(5x^7-3x+1)$。
\item \basic 求 $\deg\big((x^2+1)(x^3-2)\big)$。
\item \basic 举出两个二次多项式，其和是一次多项式。
\item \basic 求 $p(2)$，其中 $p(x)=x^3-4x+1$。
\item \basic 求 $p(-1)$，其中 $p(x)=x^4+x^2+1$。
\item \basic 判断 $x-1$ 是否整除 $x^3-1$。
\item \basic 判断 $x+2$ 是否整除 $x^2+3x+2$。
\item \basic 用带余除法把 $x^3+2x+1$ 除以 $x+1$。
\item \basic 用带余除法把 $x^4-1$ 除以 $x^2+1$。
\item \basic 求 $\gcd(x^2-1,x^2+x)$ 的首一形式。
\item \basic 求 $\gcd(x^3-x,x^2-x)$ 的首一形式。
\item \basic 证明若 $p(0)=0$，则 $x\mid p(x)$。
\item \basic 求多项式 $(x-3)^2(x+1)$ 的全部零点及其重数。
\item \basic 证明 $x^2+1$ 在 $\R[x]$ 中不可约。
\item \basic 证明 $x^2-2$ 在 $\Q[x]$ 中不可约。
\item \basic 将 $x^2-5x+6$ 在 $\R[x]$ 中分解为一次因子。
\item \basic 将 $x^4-1$ 在 $\C[x]$ 中分解为一次因子。
\end{enumerate}

\subsection*{进阶题}
\begin{enumerate}[resume, label=\textbf{\arabic*.}, leftmargin=*, itemsep=0.5em]
\item \medium 证明若 $p,q\neq 0$，则 $\deg(pq)=\deg p+\deg q$。
\item \medium 设 $\deg p\neq \deg q$，证明 $\deg(p+q)=\max\{\deg p,\deg q\}$。
\item \medium 设 $p\mid q$ 且 $q\neq 0$，证明 $\deg p\le \deg q$。
\item \medium 求 $x^4+x^3+1$ 除以 $x^2+x+1$ 的商与余式。
\item \medium 求 $x^5-1$ 除以 $x-1$ 的商。
\item \medium 证明带余除法中的商与余式唯一。
\item \medium 用欧几里得算法求 $\gcd(x^3-1,x^2-1)$。
\item \medium 用欧几里得算法求 $\gcd(x^4-1,x^3-1)$。
\item \medium 证明若 $\gcd(p,q)=1$，则存在 $a,b\in\F[x]$ 使 $ap+bq=1$。
\item \medium 设 $\gcd(p,q)=1$ 且 $p\mid qr$，证明 $p\mid r$。
\item \medium 证明若 $p(\lambda)=p(\mu)=0$ 且 $\lambda\neq \mu$，则 $(x-\lambda)(x-\mu)\mid p$。
\item \medium 证明非零多项式不可能在 $\F$ 中有多于其次数个互异零点。
\item \medium 设 $p\in\F[x]$ 的次数为 $n$，且有 $n+1$ 个互异零点，证明 $p=0$。
\item \medium 若 $p\in\R[x]$ 是奇次多项式，说明它未必在本章范围内一定有实零点；并解释为什么这不影响我们的理论。
\item \medium 证明 $\C[x]$ 中每个不可约多项式都是一次的。
\item \medium 证明 $\R[x]$ 中每个不可约多项式次数至多为 $2$。
\item \medium 证明若 $p\in\R[x]$ 且 $z\notin\R$ 是零点，则 $\overline z$ 也是零点。
\item \medium 把 $x^4+1$ 在 $\R[x]$ 中分解为两个二次因子。
\item \medium 把 $x^4+4$ 在 $\R[x]$ 中分解为实系数不可约因子。
\item \medium 设 $T\in\End(V)$，$p(x)=1+2x+x^2$，写出 $p(T)$。
\item \medium 证明对任意 $p,q$ 与 $T$，有 $(p+q)(T)=p(T)+q(T)$。
\item \medium 证明对任意 $p,q$ 与 $T$，有 $(pq)(T)=p(T)q(T)$。
\item \medium 证明 $p(T)q(T)=q(T)p(T)$。
\item \medium 设 $p(x)=(x-2)(x+1)$，把 $p(T)$ 展开成 $T$ 的幂的线性组合。
\item \medium 设 $p(x)=x^2-3x+2$。若 $p(T)=0$，证明 $(T-\Id)(T-2\Id)=0$。
\end{enumerate}

\subsection*{挑战题}
\begin{enumerate}[resume, label=\textbf{\arabic*.}, leftmargin=*, itemsep=0.5em]
\item \hard 设 $p,q\in\F[x]$ 且 $q\neq 0$。证明存在唯一的 $s,r$ 使 $p=qs+r$ 且 $\deg r<\deg q$ 或 $r=0$。
\item \hard 设 $d=\gcd(p,q)$。证明存在 $a,b\in\F[x]$ 使 $ap+bq=d$。
\item \hard 设 $\lambda_1,\dots,\lambda_n$ 两两不同，证明存在唯一次数小于 $n$ 的多项式 $p$ 满足 $p(\lambda_j)=c_j$。
\item \hard 设 $p\in\F[x]$，$\deg p=n$，且 $p$ 有 $n$ 个互异零点 $\lambda_1,\dots,\lambda_n$。证明 $p$ 必是某个非零常数乘以 $\prod_{j=1}^n(x-\lambda_j)$。
\item \hard 设 $p\in\R[x]$。证明 $p$ 在 $\R[x]$ 中不可约当且仅当 $\deg p=1$，或 $\deg p=2$ 且 $p$ 没有实零点。
\item \hard 设 $T\in\End(V)$，$p,q\in\F[x]$ 且 $p=q s$。证明 $q(T)$ 的像空间被 $p(T)$ 的像空间包含的说法一般是错的，并给出正确的可证结论。
\item \hard 设 $T\in\End(V)$，$p,q\in\F[x]$ 互素。证明存在 $a,b\in\F[x]$ 使 $a(T)p(T)+b(T)q(T)=\Id$。
\item \hard 设 $T\in\End(V)$ 且 $(T-2\Id)(T+\Id)=0$。证明对每个 $v\in V$，向量 $v$ 都可写成某个满足 $Tu=2u$ 的向量与某个满足 $Tw=-w$ 的向量之和。
\item \hard 设 $p(T)=0$，其中 $p\in\C[x]$ 非零。证明若 $p$ 没有重因子并分解为互异一次因子的乘积，则
$V$ 是这些线性因子对应核空间的直和。
\item \hard 设 $T\in\End(V)$，$p(T)=0$，其中 $p\in\R[x]$ 分解为一次与不可约二次因子的乘积。说明这为何预示着实情形中会出现一维与二维不变子空间。
\end{enumerate}

\section*{习题解答}
\begin{enumerate}[label=\textbf{\arabic*.}, leftmargin=*, itemsep=1em]
\item \textbf{解答.} 次数为 $3$。系数依次为常数项 $3$、$x$ 项 $-2$、$x^2$ 项 $0$、$x^3$ 项 $1$。
\item \textbf{解答.} 逐项相加得 $3-2x+4x^2$。
\item \textbf{解答.} 展开得 $1-x+x^2+x-x^2+x^3=1+x^3$。
\item \textbf{解答.} 最高次非零项是 $5x^7$，故次数为 $7$。
\item \textbf{解答.} 次数为 $2+3=5$。
\item \textbf{解答.} 例如 $(x^2+1)+(-x^2+x)=x+1$。
\item \textbf{解答.} $p(2)=8-8+1=1$。
\item \textbf{解答.} $p(-1)=1+1+1=3$。
\item \textbf{解答.} 因为 $1^3-1=0$，由因子定理知 $x-1\mid x^3-1$。
\item \textbf{解答.} 因为 $(-2)^2+3(-2)+2=0$，故 $x+2\mid x^2+3x+2$。
\item \textbf{解答.} 计算得 $x^3+2x+1=(x+1)(x^2-x+3)-2$，故商为 $x^2-x+3$，余式为 $-2$。
\item \textbf{解答.} 因为 $x^4-1=(x^2+1)(x^2-1)$，故商为 $x^2-1$，余式为 $0$。
\item \textbf{解答.} $x^2-1=(x-1)(x+1)$，$x^2+x=x(x+1)$，故首一最大公因子为 $x+1$。
\item \textbf{解答.} $x^3-x=x(x-1)(x+1)$，$x^2-x=x(x-1)$，故首一最大公因子为 $x(x-1)=x^2-x$。
\item \textbf{解答.} $p(0)=0$ 表示 $0$ 是零点。由因子定理，$x-0=x$ 整除 $p$。
\item \textbf{解答.} 零点为 $3$ 与 $-1$。其中 $3$ 的重数为 $2$，$-1$ 的重数为 $1$。
\item \textbf{解答.} 若 $x^2+1$ 在 $\R[x]$ 中可约，则必有实零点。可是在 $\R$ 中 $x^2+1>0$，故无实零点，所以不可约。
\item \textbf{解答.} 若在 $\Q[x]$ 可约，则有有理零点 $r$，满足 $r^2=2$。这迫使 $r=\pm\sqrt2\notin\Q$，矛盾。
\item \textbf{解答.} $x^2-5x+6=(x-2)(x-3)$。
\item \textbf{解答.} $x^4-1=(x-1)(x+1)(x-i)(x+i)$。
\item \textbf{解答.} 设首项分别为 $a_mx^m,b_nx^n$。则 $pq$ 的最高次项为 $a_mb_nx^{m+n}$，系数非零，所以次数相加。
\item \textbf{解答.} 不妨设 $\deg p=m>n=\deg q$。则 $p+q$ 的 $x^m$ 项仍为 $p$ 的首项，故次数为 $m$。
\item \textbf{解答.} 若 $q=ps$，且 $s\neq 0$，则 $\deg q=\deg p+\deg s\ge \deg p$。若 $s=0$，则 $q=0$ 与题设矛盾。
\item \textbf{解答.} 长除可得
$x^4+x^3+1=(x^2+x+1)(x^2)+(-x^2+1)=(x^2+x+1)(x^2-1)+(x+2)$。
故商为 $x^2-1$，余式为 $x+2$。
\item \textbf{解答.} 利用等比和公式，
$x^5-1=(x-1)(x^4+x^3+x^2+x+1)$。
故商为 $x^4+x^3+x^2+x+1$。
\item \textbf{解答.} 若 $p=qs_1+r_1=qs_2+r_2$，则 $q(s_1-s_2)=r_2-r_1$。若 $s_1\neq s_2$，左边次数至少为 $\deg q$，右边次数却小于 $\deg q$，矛盾，故唯一。
\item \textbf{解答.} 因为 $x^3-1=(x-1)(x^2+x+1)$，$x^2-1=(x-1)(x+1)$，故最大公因子为 $x-1$。
\item \textbf{解答.} $x^4-1=(x-1)(x^3+x^2+x+1)$，而 $x^3-1=(x-1)(x^2+x+1)$，故最大公因子为 $x-1$。
\item \textbf{解答.} 设欧几里得算法最后一个非零余式为 $d$。回代即可把 $d$ 写成 $ap+bq$。若 $\gcd(p,q)=1$，便得 $ap+bq=1$。
\item \textbf{解答.} 由上一题存在 $a,b$ 使 $ap+bq=1$。乘以 $r$ 后得 $apr+bqr=r$。右边两项都被 $p$ 整除，所以 $p\mid r$。
\item \textbf{解答.} 由因子定理，$x-\lambda$ 与 $x-\mu$ 都整除 $p$。由于它们互素，故其乘积也整除 $p$。
\item \textbf{解答.} 若有 $k$ 个互异零点，则相应线性因子乘积整除 $p$，故 $k\le\deg p$。
\item \textbf{解答.} 非零次数 $n$ 多项式至多有 $n$ 个互异零点，所以若有 $n+1$ 个互异零点，只能是零多项式。
\item \textbf{解答.} 本章未证明“奇次实系数多项式必有实零点”，因为那需要实分析中的连续性理论。我们只需要复数域上的分解与实数域上的共轭配对，这已经足够服务后续线性代数。
\item \textbf{解答.} 设 $p\in\C[x]$ 不可约且次数至少 $2$。代数基本定理给出一个复零点 $\lambda$，于是 $(x-\lambda)\mid p$，故 $p$ 可约，矛盾。
\item \textbf{解答.} 任意实系数多项式在 $\C[x]$ 中可分解为一次因子。把非实根按共轭配对，可得在 $\R[x]$ 中分解为一次或二次因子，因此不可约者次数至多为 $2$。
\item \textbf{解答.} 取共轭即可：$0=\overline{p(z)}=p(\overline z)$，因为所有系数都实。
\item \textbf{解答.} 复根为 $e^{i\pi/4},e^{3i\pi/4},e^{5i\pi/4},e^{7i\pi/4}$。配对后得
$x^4+1=(x^2+\sqrt2 x+1)(x^2-\sqrt2 x+1)$。
\item \textbf{解答.} Sophie Germain 分解：
$x^4+4=(x^2-2x+2)(x^2+2x+2)$。
两因子判别式均为 $-4$，在 $\R[x]$ 中不可约。
\item \textbf{解答.} $p(T)=\Id+2T+T^2$。
\item \textbf{解答.} 设 $p=\sum a_jx^j$，$q=\sum b_jx^j$。则 $(p+q)(T)=\sum (a_j+b_j)T^j=\sum a_jT^j+\sum b_jT^j=p(T)+q(T)$。
\item \textbf{解答.} 设 $p=\sum a_jx^j$，$q=\sum b_kx^k$。则
$p(T)q(T)=\sum_{j,k}a_jb_kT^{j+k}$，这正是按多项式乘法得到的 $(pq)(T)$。
\item \textbf{解答.} 由上一题，$p(T)q(T)=(pq)(T)=(qp)(T)=q(T)p(T)$。
\item \textbf{解答.} $(x-2)(x+1)=x^2-x-2$，故 $p(T)=T^2-T-2\Id$。
\item \textbf{解答.} 因 $p(x)=x^2-3x+2=(x-1)(x-2)$，所以
$p(T)=T^2-3T+2\Id=(T-\Id)(T-2\Id)$。若 $p(T)=0$，便得结论。
\item \textbf{解答.} 见正文带余除法定理的证明：对 $\deg p$ 归纳，逐步消去最高次项得存在性；由次数比较得唯一性。
\item \textbf{解答.} 设欧几里得算法最后一个非零余式为 $d$，则 $d$ 整除 $p,q$。回代把 $d$ 表成 $ap+bq$。任何公共因子都整除右边，所以整除 $d$，故 $d=\gcd(p,q)$。
\item \textbf{解答.} 这是拉格朗日插值。令
\[
L_j(x)=\prod_{k\ne j}\frac{x-\lambda_k}{\lambda_j-\lambda_k},
\]
则 $L_j(\lambda_i)=\delta_{ij}$。故
\[
p(x)=\sum_{j=1}^n c_jL_j(x)
\]
满足条件，且次数小于 $n$。若有两个这样的多项式，其差有 $n$ 个互异零点而次数小于 $n$，只能为零，因此唯一。
\item \textbf{解答.} 由因子定理，$\prod_{j=1}^n(x-\lambda_j)$ 整除 $p$。两者次数都为 $n$，故商为非零常数。
\item \textbf{解答.} “若”已在正文证明。反过来，若 $p\in\R[x]$ 不可约，则次数不超过 $2$。若次数为 $2$ 且有实零点，则可分解为两个一次因子，矛盾；故必须无实零点。
\item \textbf{解答.} 一般并不存在 $\range q(T)\subseteq \range p(T)$。例如取 $T=\Id$，$p=x-1$，$q=1$，则 $p=q(x-1)$，但 $p(T)=0$，$q(T)=\Id$，显然 $\range q(T)=V\not\subseteq\{0\}=\range p(T)$。正确结论是：若 $p=qs$，则 $p(T)=q(T)s(T)=s(T)q(T)$，因而
$\nul q(T)\subseteq \nul p(T)$，且 $\range p(T)\subseteq \range q(T)$。
\item \textbf{解答.} 由互素性存在 $a,b\in\F[x]$ 使 $ap+bq=1$。把 $x$ 替换为 $T$，利用多项式作用于算子的乘法性质，得
$a(T)p(T)+b(T)q(T)=\Id$。
\item \textbf{解答.} 由上一题对 $p=x-2$、$q=x+1$ 的应用，存在多项式 $a,b$ 使
$a(T)(T-2\Id)+b(T)(T+\Id)=\Id$。
对任意 $v$，令
$u=b(T)(T+\Id)v$，$w=a(T)(T-2\Id)v$，则 $v=u+w$。
再由 $(T-2\Id)(T+\Id)=0$ 以及多项式算子可交换，得
$(T-2\Id)u=0$，$(T+\Id)w=0$，故 $Tu=2u$，$Tw=-w$。
\item \textbf{解答.} 设
$p=(x-\lambda_1)\cdots(x-\lambda_m)$，各 $\lambda_j$ 互异。
这些因子两两互素。对每个 $j$，令
$q_j=\prod_{k\ne j}(x-\lambda_k)$。
由中国剩余思想或 Bezout 关系，可找到多项式 $a_j$ 使
$\sum a_jq_j=1$。代入 $T$ 得
$\sum a_j(T)q_j(T)=\Id$。
对任意 $v$，令 $v_j=a_j(T)q_j(T)v$，则 $v=\sum v_j$。
而对 $k\ne j$，$(T-\lambda_j\Id)q_j(T)$ 含有全部因子，故乘上适当多项式后为 $p(T)=0$，从而
$(T-\lambda_j\Id)v_j=0$。于是 $v_j\in \nul(T-\lambda_j\Id)$。
由不同特征值对应特征向量线性无关，可得和为直和。
\item \textbf{解答.} 一次因子 $(x-\lambda)$ 对应关系 $(T-\lambda\Id)u=0$，即一维特征方向。不可约二次因子 $x^2+bx+c$ 对应
$(T^2+bT+c\Id)u=0$。
对由某向量 $u$ 张成的循环子空间 $\spn(u,Tu)$ 而言，这给出最多二维的闭合关系，因此自然会出现二维不变子空间。后面研究实算子时，这正是实上三角化失败后仍能得到的基本块。
\end{enumerate}
